\documentclass[11pt]{article}
\usepackage[utf8]{inputenc}
\usepackage[T1]{fontenc}
\usepackage{amsmath}
\usepackage{multicol}
\usepackage{geometry}
\usepackage{tikz}
\usetikzlibrary{shapes.geometric, arrows.meta}
\usepackage{enumitem}
\usepackage{xcolor}
\usepackage{titlesec}

% Configurações de layout
\geometry{a4paper, left=1cm, right=1cm, top=0.5cm, bottom=1.2cm}
\setlength{\columnseprule}{0.4pt}
\setlength{\baselineskip}{1.0\baselineskip}

% Cores para títulos
\titleformat{\section}{\normalfont\Large\bfseries\color{blue}}{\thesection}{1em}{}
\titleformat{\subsection}{\normalfont\large\bfseries\color{red}}{\thesubsection}{1em}{}
\titleformat{\subsubsection}{\normalfont\normalsize\bfseries\color{black}}{\thesubsubsection}{1em}{}

\title{\textcolor{blue}{Função do 2º Grau: Parábolas e Aplicações}}
\author{Professor: Jefferson}
\date{}

\begin{document}

\maketitle
\vspace{-1cm}

\begin{center}
\large{Nome: \underline{\hspace{8cm}} \quad Turma: \underline{\hspace{3cm}}}
\end{center}

\begin{multicols}{2}

\section*{1. Conceito}
Uma função do 2º grau (ou quadrática) é expressa por:
\[
f(x) = ax^2 + bx + c \quad \text{ou} \quad y = ax^2 + bx + c
\]
onde:
\begin{itemize}
    \item $a$, $b$, $c$ são coeficientes reais ($a \neq 0$)
    \item $x$ é a variável independente
    \item O gráfico é sempre uma \textbf{parábola}
\end{itemize}

\subsection*{Exemplos}
\begin{itemize}
    \item $f(x) = x^2 - 4x + 3$ \quad ($a=1$, $b=-4$, $c=3$)
    \item $y = -2x^2 + 8x$ \quad ($a=-2$, $b=8$, $c=0$)
    \item $f(x) = \frac{1}{2}x^2 + 1$ \quad ($a=\frac{1}{2}$, $b=0$, $c=1$)
\end{itemize}

\section*{2. Gráfico: Parábola}
Características principais:
\begin{itemize}
    \item \textbf{Concavidade}: Para cima ($a>0$) ou para baixo ($a<0$)
    \item \textbf{Eixo de simetria}: Linha vertical que divide a parábola em duas partes simétricas
\end{itemize}

\section*{3. Zeros da Função (Raízes)}
Soluções da equação $ax^2 + bx + c = 0$:
\[
x = \frac{-b \pm \sqrt{b^2 - 4ac}}{2a}
\]
\begin{itemize}
    \item $\Delta > 0$: Duas raízes reais distintas
    \item $\Delta = 0$: Uma raiz real dupla
    \item $\Delta < 0$: Nenhuma raiz real
\end{itemize}

\section*{4. Exercícios}

\subsection*{A. Básicos (1-10)}
\begin{enumerate}
    \item Dada $f(x) = x^2 - 5x + 6$, determine:
    \begin{enumerate}[label=\alph*)]
        \item Os coeficientes $a$, $b$, $c$
        \item As raízes da função
        \item O valor de $f(0)$
        \item O valor de $f(1)$
        \item O valor de $f(2)$
    \end{enumerate}
    
    \item Classifique a concavidade:
    \begin{enumerate}[label=\alph*)]
        \item $y = 3x^2 - 2x + 1$
        \item $f(x) = -x^2 + 4$
        \item $y = \frac{1}{2}x^2 - 3x$
        \item $f(x) = -3x^2 + x - 2$
        \item $y = 0.5x^2 + 1$
    \end{enumerate}
    
    \item Calcule $\Delta$ para:
    \begin{enumerate}[label=\alph*)]
        \item $x^2 - 6x + 9 = 0$
        \item $2x^2 + x - 3 = 0$
        \item $x^2 + 4x + 5 = 0$
        \item $-3x^2 + 7x - 2 = 0$
        \item $5x^2 - 10x + 5 = 0$
    \end{enumerate}
    
    \item Resolva as equações:
    \begin{enumerate}[label=\alph*)]
        \item $x^2 - 5x + 6 = 0$
        \item $2x^2 + 3x - 2 = 0$
        \item $x^2 - 4x + 4 = 0$
        \item $3x^2 - x - 2 = 0$
        \item $-x^2 + 2x - 1 = 0$
    \end{enumerate}
\end{enumerate}

\subsection*{B. Intermediários (11-20)}
\begin{enumerate}\setcounter{enumi}{4}
    \item Para cada função, determine:
    \begin{enumerate}[label=\alph*)]
        \item As raízes de $f(x) = x^2 - 4x + 3$
        \item As raízes de $f(x) = -x^2 + 2x + 3$
        \item O valor de $f(2)$ para $f(x) = 2x^2 - 3x + 1$
        \item O valor de $f(-1)$ para $f(x) = x^2 + x - 1$
        \item O valor de $f(0.5)$ para $f(x) = 4x^2 - 4x + 1$
    \end{enumerate}
    
    \item Aplicações:
    \begin{enumerate}[label=\alph*)]
        \item Uma bola é lançada com $h(t) = -5t^2 + 20t$. Em que instantes ela atinge 15m de altura?
        \item A área de um retângulo é dada por $A(x) = -x^2 + 10x$. Para que valores de x a área é máxima?
        \item O lucro de uma empresa é dado por $L(x) = -x^2 + 50x - 400$. Qual o lucro máximo?
        \item A trajetória de um projétil é dada por $y = -x^2 + 6x$. Qual a altura máxima atingida?
        \item O custo de produção é dado por $C(x) = 0.1x^2 - 2x + 10$. Para qual quantidade o custo é mínimo?
    \end{enumerate}
    
    \item Determine $m$ para que:
    \begin{enumerate}[label=\alph*)]
        \item $f(x) = (m-1)x^2 + 2x - 3$ tenha concavidade para cima
        \item $f(x) = (2-m)x^2 - 4x + 1$ tenha concavidade para baixo
        \item $f(x) = mx^2 + 3x - 2$ tenha duas raízes reais distintas
        \item $f(x) = x^2 + mx + 9$ tenha uma única raiz real
        \item $f(x) = -x^2 + mx - m$ não tenha raízes reais
    \end{enumerate}
\end{enumerate}

\subsection*{C. Avançados (21-30)}
\begin{enumerate}\setcounter{enumi}{7}
    \item Análise gráfica:
    \begin{enumerate}[label=\alph*)]
        \item Para $f(x) = x^2 - 4x + k$, determine $k$ para que o gráfico tangencie o eixo $x$
        \item Para $f(x) = -x^2 + 2x + c$, determine $c$ para que a função tenha raízes reais
        \item Qual deve ser o valor de $p$ para que $f(x) = x^2 + px + 16$ tenha mínimo em $x = 4$?
        \item Determine $q$ para que $f(x) = x^2 + qx + 25$ tenha uma raiz dupla
        \item Para que valores de $r$ a função $f(x) = rx^2 + 4x + 1$ não tem raízes reais?
    \end{enumerate}
    
    \item Problemas complexos:
    \begin{enumerate}[label=\alph*)]
        \item Um fazendeiro quer cercar um galinheiro retangular usando um muro como um dos lados. Se ele tem 40m de cerca, qual a área máxima possível?
        \item Uma empresa estima que o custo $C(x) = 0.1x^2 - 10x + 1000$ e a receita $R(x) = 50x$. Determine quando não há lucro nem prejuízo.
        \item Num terreno retangular de 100m de perímetro, determine as dimensões para área máxima.
        \item A soma de dois números é 20. Determine-os para que o produto seja máximo.
        \item Um jardim retangular tem área de 200m². Quais suas dimensões para que o perímetro seja mínimo?
    \end{enumerate}
\end{enumerate}

\end{multicols}

\end{document}
